% ============================================================
%  Compte rendu des entretiens
%  Université Ibn Tofaïl — ENSA Kénitra
%  Module : Intelligence Artificielle — RST S6
%  Année universitaire : 2025 / 2026
% ============================================================
\documentclass[12pt,a4paper]{article}

\usepackage[utf8]{inputenc}
\usepackage[T1]{fontenc}
\usepackage[english]{babel}
\usepackage{geometry}
\usepackage{booktabs}
\usepackage{array}
\usepackage{enumitem}
\usepackage{fancyhdr}
\usepackage{titlesec}
\usepackage{hyperref}
\usepackage{parskip}
\usepackage{amssymb}
\usepackage{setspace}

\geometry{margin=2.5cm, top=3cm, bottom=3cm}

\newcolumntype{B}[1]{>{\bfseries\raggedright\arraybackslash}p{#1}}
\newcolumntype{L}[1]{>{\raggedright\arraybackslash}p{#1}}

% --- Section formatting ---
\titleformat{\section}
  {\normalfont\large\bfseries}{\thesection.}{0.6em}{}
\titlespacing{\section}{0pt}{16pt}{6pt}

\titleformat{\subsection}
  {\normalfont\normalsize\bfseries}{\thesubsection.}{0.5em}{}
\titlespacing{\subsection}{0pt}{10pt}{4pt}

% --- Header/Footer ---
\pagestyle{fancy}
\fancyhf{}
\fancyhead[L]{\small Compte rendu des entretiens --- ENSA K\'{e}nitra}
\fancyhead[R]{\small 2025 / 2026}
\fancyfoot[C]{\small Page \thepage{} / 18}
\renewcommand{\headrulewidth}{0.4pt}

\hypersetup{colorlinks=false, hidelinks}
\setlength{\parskip}{6pt}
\setlength{\parindent}{0pt}
\setlength{\tabcolsep}{6pt}

% Guillemets
\newcommand{\og}{\guillemotleft\,}
\newcommand{\fg}{\,\guillemotright}

\begin{document}

% ============================================================
%  TITLE PAGE
% ============================================================
\begin{titlepage}
  \centering
  \vspace*{1.5cm}
  {\large Universit\'{e} Ibn Tofa\"{i}l}\\[0.3em]
  {\normalsize \'{E}cole Nationale des Sciences Appliqu\'{e}es --- K\'{e}nitra}
  \vspace{1.5cm}

  \rule{\textwidth}{0.8pt}\\[0.6em]
  {\LARGE\bfseries Compte rendu des entretiens}\\[0.5em]
  {\large Recueil de t\'{e}moignages et retours d'exp\'{e}rience professionnels}\\[0.4em]
  \rule{\textwidth}{0.4pt}

  \vspace{1.5cm}
  \begin{tabular}{@{}ll@{}}
    \textbf{Module :}          & Intelligence Artificielle \\[5pt]
    \textbf{Fili\`{e}re :}    & R\'{e}seaux et Syst\`{e}mes de T\'{e}l\'{e}communications --- RST --- S6 \\[5pt]
    \textbf{Encadrant :}       & Aniss Moumen \\[5pt]
    \textbf{Membres du groupe :} & Nafia Khadija, El Kartouti Anas, Lekhbioui Adil, \\
                               & Moulay Anas, Hammoubel El Yazid \\[5pt]
    \textbf{Ann\'{e}e universitaire :} & 2025 / 2026 \\[5pt]
    \textbf{Date de remise :}  & Juin 2026 \\
  \end{tabular}

  \vfill
  {\small ENSA K\'{e}nitra --- R\'{e}seaux et Syst\`{e}mes de T\'{e}l\'{e}communications}
\end{titlepage}

% ============================================================
%  TABLE DES MATIÈRES
% ============================================================
\tableofcontents
\clearpage

% ============================================================
%  INTRODUCTION
% ============================================================
\section*{Introduction}
\addcontentsline{toc}{section}{Introduction}

Dans le cadre de notre projet de fin de module d'Intelligence Artificielle, nous nous sommes
int\'{e}ress\'{e}s \`{a} l'utilisation du machine learning pour la d\'{e}tection des cyberattaques, en
particulier les attaques bas\'{e}es sur le scan de ports. Afin d'enrichir notre approche th\'{e}orique
et de confronter nos travaux \`{a} la r\'{e}alit\'{e} du terrain, nous avons jug\'{e} utile de rencontrer des
experts du domaine de la cybers\'{e}curit\'{e} et des r\'{e}seaux.

Ces entretiens avaient pour objectif de mieux comprendre comment les professionnels font
face \`{a} ce type de menaces dans leur environnement de travail, quels outils et m\'{e}thodes ils
utilisent, et dans quelle mesure l'intelligence artificielle est aujourd'hui int\'{e}gr\'{e}e dans les
strat\'{e}gies de d\'{e}fense r\'{e}seau.

Ce rapport regroupe l'ensemble des cinq entretiens r\'{e}alis\'{e}s avec des ing\'{e}nieurs sp\'{e}cialis\'{e}s,
suivis d'une synth\`{e}se des principaux enseignements tir\'{e}s de ces \'{e}changes.

\clearpage

% ============================================================
%  ENTRETIEN 1 — HICHAM DACHOUR
% ============================================================
\section{Entretien 1 --- Hicham Dachour}

\subsection*{Informations sur l'interlocuteur}

\begin{tabular}{@{}B{4cm} L{10.5cm}@{}}
  \hline
  Nom et pr\'{e}nom      & Hicham Dachour \\
  \hline
  Poste / Fonction   & Cybersecurity \& Network Engineer $|$ CCNP ENCOR $|$ Fortinet NSE7
                       Certified (Enterprise Firewall, SD-WAN) $|$ Advanced Infrastructure
                       Security : WAF \& Load Balancing $|$ Research Scholar : AI/ML
                       Integration in Cyber Resilience \\
  \hline
  Entreprise         & Tuntrust \\
  \hline
  Secteur            & Cybers\'{e}curit\'{e} \\
  \hline
  Date de l'entretien & 26 avril 2026 \\
  \hline
  Lieu / Format      & En ligne --- Google Meet \\
  \hline
\end{tabular}

\subsection*{Questions et r\'{e}ponses}

\textbf{Q1. Pr\'{e}sentation et prise de contact}

Bonjour, M. Hicham. J'esp\`{e}re que vous allez bien. Je me pr\'{e}sente : je m'appelle Khadija Nafia.
Je suis \'{e}tudiante en premi\`{e}re ann\'{e}e du cycle ing\'{e}nieur \`{a} l'\'{E}cole Nationale des Sciences
Appliqu\'{e}es, sp\'{e}cialit\'{e} R\'{e}seaux et Syst\`{e}mes de T\'{e}l\'{e}communication. Dans le cadre de notre
projet de fin de module en Intelligence Artificielle, nous travaillons sur la d\'{e}tection des attaques
r\'{e}seau \`{a} partir de techniques de machine learning. L'objectif de cet entretien est de recueillir
votre exp\'{e}rience et votre point de vue sur la probl\'{e}matique de la s\'{e}curit\'{e} r\'{e}seau, ainsi que
sur l'apport potentiel de l'intelligence artificielle dans ce domaine. Je vous remercie une nouvelle
fois d'avoir accept\'{e} de participer \`{a} cet \'{e}change et pour votre disponibilit\'{e}. Je vous informe
\'{e}galement que cet entretien sera enregistr\'{e} dans le cadre de notre projet. \^{E}tes-vous \`{a} l'aise
avec cela ?

\medskip
\textit{R\'{e}ponse}

Oui, d'accord.

\bigskip
\textbf{Q2. Parcours et r\^{o}le actuel}

Merci beaucoup. Je m'appelle Hicham Dachraoui. Je suis ing\'{e}nieur en cybers\'{e}curit\'{e} avec plus
de quatre ans d'exp\'{e}rience dans la conception, l'optimisation et la s\'{e}curisation d'infrastructures
d'entreprise. Actuellement en poste d'ing\'{e}nieur en cybers\'{e}curit\'{e}, je m\`{e}ne \'{e}galement des
recherches sur l'utilisation de l'IA dans ce domaine. Je supervise des solutions de s\'{e}curit\'{e}
avanc\'{e}es, notamment les pare-feux de nouvelle g\'{e}n\'{e}ration (Next-Generation Firewall) et les
pare-feux applicatifs web (WAF) destin\'{e}s \`{a} prot\'{e}ger les applications publi\'{e}es, ainsi que des
solutions SIEM (Security Information and Event Management). Les architectures que je g\`{e}re sont
principalement hybrides, combinant des environnements LAN d'entreprise, des infrastructures
cloud et des connexions s\'{e}curis\'{e}es via des tunnels VPN point \`{a} point ou site \`{a} site.
J'interviens \'{e}galement sur des architectures Zero Trust et le d\'{e}ploiement de solutions de centre
de donn\'{e}es telles que Cisco ACI et Cisco Nexus.

\bigskip
\textbf{Q3. Principaux d\'{e}fis quotidiens}

Les d\'{e}fis sont nombreux : la r\'{e}duction des menaces avanc\'{e}es telles que les attaques
zero-day (des attaques pour lesquelles aucun correctif n'existe encore), la gestion du volume
de journaux et d'alertes, la complexit\'{e} croissante face aux attaques DDoS (Distributed Denial
of Service), la conformit\'{e} r\'{e}glementaire avec des politiques de s\'{e}curit\'{e} fond\'{e}es sur la norme
ISO 27001, et la s\'{e}curisation des acc\`{e}s privil\'{e}gi\'{e}s via des m\'{e}canismes comme le MFA
(Multi-Factor Authentication).

\bigskip
\textbf{Q4. Comportements inhabituels observ\'{e}s}

Lors des op\'{e}rations de surveillance, nous observons fr\'{e}quemment des scans actifs de ports,
notamment par l'utilisation de l'outil Nmap, visant une plage d'adresses IP ou un h\^{o}te particulier
h\'{e}bergeant un site web ou une application. L'attaquant tente de recenser les ports ouverts afin
de pr\'{e}parer une attaque. La r\'{e}ponse consiste \`{a} appliquer des r\`{e}gles de s\'{e}curit\'{e} via les
solutions de pare-feu de nouvelle g\'{e}n\'{e}ration, qui permettent de contrer ce type d'attaque.

\bigskip
\textbf{Q5. Causes les plus fr\'{e}quentes des anomalies}

Les causes les plus fr\'{e}quentes sont : l'activit\'{e} de reconnaissance (scans de ports,
\'{e}num\'{e}ration des services), les tentatives de force brute, le mouvement lat\'{e}ral au sein du
syst\`{e}me, l'exploitation de vuln\'{e}rabilit\'{e}s non corrig\'{e}es ou de mauvaises configurations de
s\'{e}curit\'{e}, l'exfiltration discr\`{e}te de donn\'{e}es via des canaux chiffr\'{e}s (comme le DNS tunneling),
et les attaques ciblant des syst\`{e}mes industriels ou des services tiers.

\bigskip
\textbf{Q6. Diff\'{e}rencier une panne technique d'une attaque malveillante}

Une panne technique classique se traduit par une chute brutale du trafic, une absence de flux
ou une d\'{e}faillance mat\'{e}rielle. On peut la diagnostiquer via un outil de supervision utilisant le
protocole SNMP (agentless), qui renseigne sur la sant\'{e} des \'{e}quipements : CPU, RAM, stockage
et temp\'{e}rature. Les attaques malveillantes, quant \`{a} elles, impliquent souvent des commandes
de contr\^{o}le \`{a} distance ou du phishing. La protection passe par la mise en place d'EDR
(Endpoint Detection and Response) et de SIEM pour collecter les journaux, analyser les
incidents et am\'{e}liorer continuellement les playbooks de d\'{e}tection. La sensibilisation du
personnel est \'{e}galement essentielle pour reconna\^{i}tre les tentatives de phishing.

\bigskip
\textbf{Q7. D\'{e}tection des scans de ports}

L'analyse des journaux du pare-feu permet d'identifier les tentatives de scan. La surveillance
des requ\^{e}tes DNS constitue \'{e}galement un indicateur pertinent.

\bigskip
\textbf{Q8. Outils de surveillance et de s\'{e}curisation}

La d\'{e}tection commence par la solution SIEM, qui assure la corr\'{e}lation des \'{e}v\'{e}nements. Pour
l'investigation, on peut utiliser Wireshark. La segmentation r\'{e}seau via des VLANs permet
\'{e}galement de compartimenter le r\'{e}seau et de limiter l'interconnexion entre services distincts.

\bigskip
\textbf{Q9. Limites des outils actuels}

On observe des faux positifs, des temps de r\'{e}action parfois longs susceptibles de perturber
le r\'{e}seau, ainsi que des difficult\'{e}s li\'{e}es \`{a} la complexit\'{e} d'int\'{e}gration de ces solutions.

\bigskip
\textbf{Q10. Apport du machine learning dans la d\'{e}tection des attaques}

L'IA peut aujourd'hui aider les ing\'{e}nieurs \`{a} d\'{e}tecter des menaces en entra\^{i}nant un agent
capable de d\'{e}tecter et de r\'{e}agir automatiquement --- bloquer ou autoriser un flux. L'ann\'{e}e
derni\`{e}re, j'ai r\'{e}dig\'{e} un article sur la cr\'{e}ation d'agents IA au niveau des routeurs Cisco pour
d\'{e}tecter des \'{e}v\'{e}nements BGP, notamment le BGP hijacking, un type d'attaque qui provoque
un reroutage du trafic. L'utilisation de l'IA peut donc aider les sp\'{e}cialistes \`{a} mieux prot\'{e}ger
leur infrastructure.

\bigskip
\textbf{Q11. Blocage automatique par l'IA}

Pas syst\'{e}matiquement. Un blocage automatique permanent n'est pas toujours souhaitable.
On peut envisager un blocage temporaire (une \`{a} deux heures, par exemple en dehors des
heures ouver\'{e}es), qui peut \^{e}tre progressivement allong\'{e} si les m\^{e}mes adresses IP
renouvellent leurs tentatives, puisque les attaquants changent souvent d'adresse IP.

\bigskip
\textbf{Q12. Fonctionnalit\'{e}s essentielles d'une solution IA}

Une solution IA efficace devrait passer par une phase d'entra\^{i}nement, puis de test, avant la
mise en production. Elle pourrait inclure un tableau de bord unifi\'{e} avec la g\'{e}olocalisation des
IP, un syst\`{e}me d'alertes par e-mail ou SMS, l'int\'{e}gration via des API et la gestion des
politiques de s\'{e}curit\'{e}.

\clearpage

% ============================================================
%  ENTRETIEN 2 — SOUKAINA NAFIA
% ============================================================
\section{Entretien 2 --- Soukaina Nafia}

\subsection*{Informations sur l'interlocuteur}

\begin{tabular}{@{}B{4cm} L{10.5cm}@{}}
  \hline
  Nom et pr\'{e}nom      & Soukaina Nafia \\
  \hline
  Poste / Fonction   & Consultante en cybers\'{e}curit\'{e} chez CY-CLIC $|$ Mentore en
                       cybers\'{e}curit\'{e} chez Connected Canadians $|$ \'{E}tudiante en
                       cybers\'{e}curit\'{e} \`{a} Polytechnique Montr\'{e}al \\
  \hline
  Entreprise         & CY-CLIC / SECOM \\
  \hline
  Secteur            & Cybers\'{e}curit\'{e} \\
  \hline
  Date de l'entretien & 08 juin 2026 \\
  \hline
  Lieu / Format      & Microsoft Teams \\
  \hline
\end{tabular}

\subsection*{Questions et r\'{e}ponses}

\textbf{Q1. Pr\'{e}sentation et prise de contact}

Bonjour, je suis ravie de vous accueillir aujourd'hui. Je m'appelle Khadija et nous allons avoir
un \'{e}change tr\`{e}s int\'{e}ressant autour de l'intervention de l'intelligence artificielle dans la
d\'{e}tection du trafic malveillant. Pourriez-vous nous parler bri\`{e}vement de votre parcours et de
votre r\^{o}le actuel dans la gestion des r\'{e}seaux ?

\medskip
\textit{R\'{e}ponse}

Bonjour, merci de m'avoir invit\'{e}e. Mon nom est Soukaina Nafia. Je suis \'{e}tudiante au
baccalaur\'{e}at en cybers\'{e}curit\'{e} \`{a} Polytechnique Montr\'{e}al. En parall\`{e}le, je travaille chez SECOM
en tant que consultante et formatrice en pr\'{e}vention de la fraude et en cybers\'{e}curit\'{e}. Mon
r\^{o}le touche principalement la sensibilisation, l'analyse des risques, les bonnes pratiques de
s\'{e}curit\'{e} et l'accompagnement d'organisations souhaitant mieux prot\'{e}ger leurs donn\'{e}es. Je
fais \'{e}galement du b\'{e}n\'{e}volat chez Connected Canadians, o\`{u} j'aide \`{a} s\'{e}curiser un petit r\'{e}seau
informatique compos\'{e} d'une vingtaine d'ordinateurs. Dans mon programme \`{a} Polytechnique,
nous travaillons beaucoup en laboratoire avec des machines virtuelles, notamment sous Linux.
Nous r\'{e}alisons des exercices de reconnaissance r\'{e}seau, d'analyse de vuln\'{e}rabilit\'{e}s et de
tests d'intrusion. Dans mon cours CR470 consacr\'{e} aux tests d'intrusion, les travaux pratiques
demandent d'utiliser des environnements virtuels comme Linux, Windows 10 et Metasploitable
2 pour effectuer de la reconnaissance, du balayage r\'{e}seau et de l'analyse de vuln\'{e}rabilit\'{e}s.

\bigskip
\textbf{Q2. Principaux d\'{e}fis quotidiens}

Dans une petite organisation, le plus grand d\'{e}fi est souvent le manque de ressources. Il n'y a
pas toujours une personne d\'{e}di\'{e}e \`{a} temps plein \`{a} la s\'{e}curit\'{e} informatique. Les principaux
d\'{e}fis sont la gestion des mots de passe, les mises \`{a} jour, le contr\^{o}le des acc\`{e}s, la protection
des postes de travail et l'encadrement des acc\`{e}s \`{a} distance par VPN. Il faut \'{e}galement veiller
\`{a} appliquer le principe du moindre privil\`{e}ge. Dans un petit r\'{e}seau, la visibilit\'{e} est aussi un
d\'{e}fi : on ne sait pas toujours quels appareils sont connect\'{e}s, quels services sont expos\'{e}s ou
si certains postes pr\'{e}sentent des comportements anormaux. C'est pourquoi l'inventaire, la
surveillance minimale et les bonnes pratiques de base sont essentiels.

\bigskip
\textbf{Q3. Comportements inhabituels observ\'{e}s}

Dans un petit r\'{e}seau, une anomalie typique peut se manifester par une latence \'{e}lev\'{e}e ou une
connexion instable pour plusieurs utilisateurs. Avant de conclure \`{a} une attaque, il faut examiner
les causes les plus simples : surcharge de la connexion Internet, probl\`{e}me de routeur, appareil
consommant une grande partie de la bande passante, mise \`{a} jour massive ou mauvaise
configuration. Dans mes laboratoires, nous apprenons \`{a} analyser ce type de situation avec une
approche structur\'{e}e : identifier les machines pr\'{e}sentes sur le r\'{e}seau, v\'{e}rifier les ports ouverts,
comprendre les services expos\'{e}s, puis d\'{e}terminer si l'activit\'{e} observ\'{e}e est normale ou
suspecte. Mon cours CR470 demande d'utiliser Nmap pour effectuer diff\'{e}rents types de scans
r\'{e}seau --- TCP, UDP et analyses d'empreintes (fingerprinting) --- afin d'identifier les syst\`{e}mes
et services pr\'{e}sents.

\bigskip
\textbf{Q4. Causes les plus fr\'{e}quentes des anomalies}

Les causes les plus courantes ne sont pas toujours malveillantes : \'{e}quipement vieillissant,
mauvaise configuration, poste non mis \`{a} jour, logiciel consommant trop de ressources ou
probl\`{e}me avec le fournisseur Internet. Mais il faut aussi envisager des causes li\'{e}es \`{a} la
s\'{e}curit\'{e} : appareil compromis, logiciel malveillant, compte VPN utilis\'{e} sans autorisation, scan
interne ou poste communiquant avec des adresses IP inhabituelles. Dans un petit organisme,
le risque vient souvent du manque de journalisation, de visibilit\'{e} et de tra\c{c}abilit\'{e}, ce qui rend
difficile la d\'{e}tection rapide d'un comportement anormal.

\bigskip
\textbf{Q5. Diff\'{e}rencier une panne technique d'une attaque malveillante}

Je commencerais par une approche globale : si tout le r\'{e}seau est lent, je v\'{e}rifierais l'\'{e}tat de
la connexion Internet, du pare-feu et des \'{e}quipements r\'{e}seau. Si seulement un poste ou un
groupe de postes est touch\'{e}, je me concentrerais sur les machines concern\'{e}es. Pour une
surcharge ou une attaque par d\'{e}ni de service (DoS/DDoS), je chercherais des signes comme
un volume anormal de connexions, un trafic entrant excessif ou des requ\^{e}tes r\'{e}p\'{e}titives vers
les m\^{e}mes services. Pour un malware interne, je rechercherais un poste communiquant avec
des destinations inhabituelles, g\'{e}n\'{e}rant beaucoup de trafic sortant ou tentant de contacter
plusieurs machines internes. Dans mes laboratoires, j'utilise des outils comme Nmap, Nessus
ou d'autres analyseurs pour identifier les services expos\'{e}s, les vuln\'{e}rabilit\'{e}s et les
comportements suspects.

\bigskip
\textbf{Q6. D\'{e}tection des scans de ports}

Un scan de ports peut se manifester par des tentatives r\'{e}p\'{e}t\'{e}es de connexion vers plusieurs
ports ou plusieurs machines en peu de temps. Dans un r\'{e}seau bien surveill\'{e}, il peut \^{e}tre
d\'{e}tect\'{e} \`{a} l'aide des journaux du pare-feu, d'un IDS, d'un IPS ou d'un outil de supervision
r\'{e}seau. Dans une petite structure, ces solutions ne sont pas toujours disponibles. Il faut donc,
au minimum, activer les journaux utiles sur les \'{e}quipements disponibles, surveiller les
connexions VPN, conserver une trace des appareils autoris\'{e}s et limiter les services expos\'{e}s.
Dans mon cours de cybers\'{e}curit\'{e}, cela inclut une partie consacr\'{e}e \`{a} la reconnaissance
OSINT avec Maltego ainsi que plusieurs scans Nmap sur des machines virtuelles.

\bigskip
\textbf{Q7. Outils de surveillance et de s\'{e}curisation}

Dans ce contexte, je m'appuierais sur les sources d'information disponibles : la liste des
appareils connect\'{e}s et autoris\'{e}s, les journaux VPN, les alertes de l'antivirus, les
\'{e}v\'{e}nements Windows et, si possible, l'observation du trafic r\'{e}seau. La premi\`{e}re \'{e}tape serait
d'identifier l'appareil ou le compte concern\'{e}. Ensuite, je recommanderais de contenir le risque
en isolant le poste suspect. Si le compte semble compromis, il faudrait imm\'{e}diatement changer
le mot de passe, r\'{e}voquer les sessions actives et v\'{e}rifier les acc\`{e}s r\'{e}cents. Je
documenterais \'{e}galement l'incident de mani\`{e}re structur\'{e}e, en gardant une approche
pragmatique adapt\'{e}e aux ressources disponibles.

\bigskip
\textbf{Q8. Limites des outils actuels}

La principale limite est le manque de temps et de ressources. Dans une petite organisation,
les employ\'{e}s ne sont pas forc\'{e}ment form\'{e}s pour interpr\'{e}ter les journaux ou reconna\^{i}tre une
anomalie r\'{e}seau. On ne dispose pas toujours d'un SIEM ni d'un IDS, ce qui augmente le
risque de manquer des signaux faibles. Les faux positifs constituent \'{e}galement une limite
importante : une activit\'{e} peut sembler suspecte alors qu'elle est caus\'{e}e par une mise \`{a} jour,
une sauvegarde ou un service l\'{e}gitime. C'est pourquoi il est essentiel de bien documenter le
comportement normal du r\'{e}seau.

\bigskip
\textbf{Q9. Apport du machine learning dans la d\'{e}tection des attaques}

Je pense que l'IA peut \^{e}tre tr\`{e}s utile pour d\'{e}tecter des anomalies plus rapidement, surtout
lorsqu'il y a beaucoup de donn\'{e}es \`{a} analyser. Un poste communiquant soudainement avec
plusieurs adresses externes, un volume de trafic anormal ou des tentatives de connexion
r\'{e}p\'{e}t\'{e}es peuvent \^{e}tre plus facilement d\'{e}tect\'{e}s. Cependant, je ne la verrais pas comme une
solution autonome d'embl\'{e}e : l'IA devrait d'abord \^{e}tre un outil d'aide \`{a} la d\'{e}cision, capable
de prioriser les alertes et d'expliquer pourquoi un comportement semble anormal. Il faut
\'{e}galement \^{e}tre attentif \`{a} la confidentialit\'{e} des donn\'{e}es collect\'{e}es. L'IA peut aider, mais elle
doit rester encadr\'{e}e et utilis\'{e}e de mani\`{e}re responsable.

\bigskip
\textbf{Q10. Blocage automatique par l'IA}

Oui, mais avec beaucoup de prudence. L'IA peut \^{e}tre utile pour recommander le blocage d'une
adresse IP ou g\'{e}n\'{e}rer une alerte prioritaire. En revanche, je serais plus r\'{e}serv\'{e}e concernant
les d\'{e}cisions automatiques --- comme le blocage d'un routeur ou la coupure d'un acc\`{e}s VPN
--- car une mauvaise d\'{e}cision pourrait bloquer tout un environnement ou interrompre un service.
La meilleure approche serait progressive : d'abord une phase de d\'{e}tection et de signalement,
puis une automatisation avec validation humaine, et enfin une autorisation uniquement pour des
sc\'{e}narios tr\`{e}s bien d\'{e}finis et \`{a} faible risque.

\bigskip
\textbf{Q11. Fonctionnalit\'{e}s essentielles d'une solution IA}

Une solution IA efficace en cybers\'{e}curit\'{e} devrait inclure : une d\'{e}tection d'anomalies en temps
r\'{e}el, un syst\`{e}me de corr\'{e}lation des \'{e}v\'{e}nements pour relier plusieurs alertes entre elles, une
priorisation intelligente des alertes, des explications claires des d\'{e}cisions de l'IA pour faciliter
la prise de d\'{e}cision, et, enfin, un contr\^{o}le humain maintenu dans la boucle pour les actions
sensibles.

\clearpage

% ============================================================
%  ENTRETIEN 3 — OUSSAMA
% ============================================================
\section{Entretien 3 --- Oussama}

\subsection*{Informations sur l'interlocuteur}

\begin{tabular}{@{}B{4cm} L{10.5cm}@{}}
  \hline
  Nom et pr\'{e}nom      & Oussama \\
  \hline
  Poste / Fonction   & \'{E}tudiant en 2$^{e}$ ann\'{e}e --- G\'{e}nie informatique --- ENSA K\'{e}nitra \\
  \hline
  Entreprise         & --- \\
  \hline
  Secteur            & Cybers\'{e}curit\'{e} (contexte acad\'{e}mique) \\
  \hline
  Date de l'entretien & 28 avril 2026 \\
  \hline
  Lieu / Format      & Microsoft Teams \\
  \hline
\end{tabular}

\subsection*{Questions et r\'{e}ponses}

\textbf{Q1. Pr\'{e}sentation et architecture r\'{e}seau utilis\'{e}e}

Je m'appelle Oussama. Je suis \'{e}tudiant en deuxi\`{e}me ann\'{e}e en g\'{e}nie informatique. Au
quotidien, j'utilise principalement le r\'{e}seau Wi-Fi de l'\'{e}cole pour acc\'{e}der aux plateformes
p\'{e}dagogiques, effectuer des recherches, travailler sur GitHub et communiquer avec mon groupe
de projet. En dehors des cours, j'utilise aussi parfois un petit environnement local avec Linux,
Docker ou des machines virtuelles pour tester des applications web ou des scripts r\'{e}seau.

\bigskip
\textbf{Q2. Principaux d\'{e}fis li\'{e}s \`{a} l'infrastructure r\'{e}seau}

Le probl\`{e}me principal est la disponibilit\'{e} et la stabilit\'{e} du r\'{e}seau. Aux heures de pointe,
lorsque de nombreux \'{e}tudiants sont connect\'{e}s, la connexion devient tr\`{e}s lente. Parfois, le
signal Wi-Fi est bon mais les pages ne chargent pas, ou les plateformes p\'{e}dagogiques
r\'{e}pondent tr\`{e}s lentement. On observe \'{e}galement des coupures temporaires, des pertes de
paquets ou une latence \'{e}lev\'{e}e.

\bigskip
\textbf{Q3. Exp\'{e}rience r\'{e}cente d'anomalie r\'{e}seau}

R\'{e}cemment, pendant un travail de groupe, nous devions envoyer un fichier et acc\'{e}der \`{a} une
plateforme en ligne, mais la connexion \'{e}tait tr\`{e}s lente. Le Wi-Fi \'{e}tait connect\'{e}, mais les
requ\^{e}tes mettaient beaucoup de temps \`{a} r\'{e}pondre, provoquant des erreurs de type time-out.
Le probl\`{e}me semblait g\'{e}n\'{e}ral, touchant plusieurs services en m\^{e}me temps, comme si une
saturation ou une anomalie affectait le trafic r\'{e}seau.

\bigskip
\textbf{Q4. Causes probables des perturbations furtives}

Plusieurs causes sont possibles : la surcharge du r\'{e}seau due \`{a} un grand nombre d'utilisateurs
connect\'{e}s, mais aussi des machines g\'{e}n\'{e}rant un trafic excessif \`{a} cause de t\'{e}l\'{e}chargements
lourds, de mises \`{a} jour automatiques ou d'un malware. Dans un contexte de s\'{e}curit\'{e}, une
machine compromise pourrait lancer des scans de ports, envoyer de nombreuses requ\^{e}tes ou
tenter de d\'{e}couvrir les services ouverts sur le r\'{e}seau. Ce type d'activit\'{e} peut \^{e}tre discret au
d\'{e}but, mais il peut n\'{e}anmoins d\'{e}grader les performances.

\bigskip
\textbf{Q5. Diff\'{e}rencier une panne d'une attaque}

\`{A} mon niveau, je ne peux pas le confirmer avec certitude, car je n'ai pas acc\`{e}s aux journaux
des routeurs ou des pare-feux. Je peux n\'{e}anmoins effectuer quelques v\'{e}rifications simples :
v\'{e}rifier si le probl\`{e}me concerne uniquement mon poste ou plusieurs \'{e}tudiants, faire un ping
vers la passerelle ou vers un site externe pour mesurer la latence et les pertes de paquets. Si
mon ordinateur est le seul affect\'{e}, je v\'{e}rifie les processus consommant de la bande passante
ou je lance une analyse antivirus.

\bigskip
\textbf{Q6. Outils et proc\'{e}dures utilis\'{e}s}

Je commence par des tests simples : v\'{e}rifier la connexion, faire un ping, tester un autre
navigateur ou un autre appareil, puis v\'{e}rifier si une application consomme beaucoup de bande
passante. Dans mes projets personnels, je peux utiliser des outils comme Wireshark ou netstat
pour observer les connexions actives. Sur le r\'{e}seau de l'\'{e}cole, je n'ai pas les droits n\'{e}cessaires
pour analyser toute l'infrastructure. Pour bloquer une menace, je peux seulement s\'{e}curiser
ma propre machine, fermer des applications suspectes, lancer un antivirus ou signaler le
probl\`{e}me au service informatique.

\bigskip
\textbf{Q7. Int\'{e}gration de l'IA pour d\'{e}tecter les attaques en temps r\'{e}el}

Je pense que c'est pertinent, car un mod\`{e}le d'IA pourrait analyser des volumes de donn\'{e}es
bien plus importants qu'un humain. Il pourrait d\'{e}tecter des patterns anormaux comme un
nombre inhabituel de connexions, des scans de ports, des tentatives r\'{e}p\'{e}t\'{e}es vers plusieurs
services ou une hausse brutale du trafic. L'avantage est sa r\'{e}activit\'{e}, surtout si le mod\`{e}le est
entra\^{i}n\'{e} sur des comportements normaux et malveillants. Il faut cependant veiller \`{a} limiter les
faux positifs.

\bigskip
\textbf{Q8. Conditions d'adoption d'une solution IA}

Oui, je serais pr\^{e}t \`{a} l'utiliser si elle am\'{e}liore la stabilit\'{e} du r\'{e}seau et permet une d\'{e}tection
rapide des anomalies. Les fonctionnalit\'{e}s importantes seraient une analyse en temps r\'{e}el,
des alertes claires et un tableau de bord simple. Mes principales craintes concernent la
confidentialit\'{e} des donn\'{e}es, le risque de faux positifs et le fait que l'IA puisse prendre des
d\'{e}cisions automatiques sans contr\^{o}le humain.

\clearpage

% ============================================================
%  ENTRETIEN 4 — OUZIZI AMINE
% ============================================================
\section{Entretien 4 --- Ouzizi Amine}

\subsection*{Informations sur l'interlocuteur}

\begin{tabular}{@{}B{4cm} L{10.5cm}@{}}
  \hline
  Nom et pr\'{e}nom      & Ouzizi Amine \\
  \hline
  Poste / Fonction   & \'{E}tudiant en 2$^{e}$ ann\'{e}e du cycle ing\'{e}nieur --- R\'{e}seaux et Syst\`{e}mes
                       de T\'{e}l\'{e}communication, option S\'{e}curit\'{e} des R\'{e}seaux et des
                       Syst\`{e}mes d'Information \\
  \hline
  Entreprise         & --- \\
  \hline
  Secteur            & Cybers\'{e}curit\'{e} (contexte acad\'{e}mique) \\
  \hline
  Date de l'entretien & 28 avril 2026 \\
  \hline
  Lieu / Format      & Microsoft Teams \\
  \hline
\end{tabular}

\subsection*{Questions et r\'{e}ponses}

\textbf{Q1. Pr\'{e}sentation et architecture r\'{e}seau utilis\'{e}e}

Je m'appelle Ouzizi Amine, \'{e}tudiant en deuxi\`{e}me ann\'{e}e du cycle ing\'{e}nieur en r\'{e}seaux et
syst\`{e}mes de t\'{e}l\'{e}communication, option s\'{e}curit\'{e} des r\'{e}seaux et des syst\`{e}mes d'information.
J'utilise le r\'{e}seau Wi-Fi de l'\'{e}cole tous les jours et je fais souvent tourner des environnements
de test locaux sur mon poste Windows/Ubuntu pour avancer sur mes projets acad\'{e}miques et
professionnels.

\bigskip
\textbf{Q2. Principaux d\'{e}fis li\'{e}s \`{a} l'infrastructure r\'{e}seau}

Le principal d\'{e}fi est l'instabilit\'{e} de la connexion qui l\^{a}che souvent en pleine session de travail
sans explication, ou des ports locaux qui deviennent soudainement inaccessibles depuis les
autres machines de mon groupe de travaux pratiques.

\bigskip
\textbf{Q3. Exp\'{e}rience r\'{e}cente d'anomalie r\'{e}seau}

La semaine derni\`{e}re, mon serveur web local plantait en boucle. En consultant les journaux de
mon syst\`{e}me, j'ai constat\'{e} des centaines de requ\^{e}tes incompl\`{e}tes sur diff\'{e}rents ports,
provenant d'adresses IP appartenant au r\'{e}seau \'{e}tudiant.

\bigskip
\textbf{Q4. Causes probables des perturbations furtives}

Souvent, ce sont des scripts automatis\'{e}s ou des malwares install\'{e}s \`{a} l'insu d'un autre \'{e}tudiant
sur son PC, qui tentent de trouver d'autres machines vuln\'{e}rables sur le r\'{e}seau local.

\bigskip
\textbf{Q5. Diff\'{e}rencier une panne d'une attaque}

Je consulte les outils de surveillance de mon poste. Si l'utilisation du CPU est normale mais
que mes ports de communication tombent les uns apr\`{e}s les autres, je sais que la perturbation
provient de l'ext\'{e}rieur.

\bigskip
\textbf{Q6. D\'{e}tection des tentatives de reconnaissance}

Sur mon syst\`{e}me, je peux les d\'{e}tecter en consultant les journaux de mon pare-feu local.
Cependant, cette v\'{e}rification manuelle est difficile \`{a} maintenir en continu.

\bigskip
\textbf{Q6b. Outils et proc\'{e}dures utilis\'{e}s}

Je consulte les connexions actives dans mon terminal et j'ajoute manuellement l'adresse IP
suspecte dans les r\`{e}gles de rejet de mon pare-feu.

\bigskip
\textbf{Q7. Int\'{e}gration de l'IA pour d\'{e}tecter les attaques en temps r\'{e}el}

Ce serait tr\`{e}s utile. L'IA pourrait comprendre mon comportement habituel en tant que
d\'{e}veloppeur sur le r\'{e}seau et bloquer uniquement ce qui constitue r\'{e}ellement un trafic de
reconnaissance malveillant.

\bigskip
\textbf{Q8. Pertinence de l'IA pour intervenir sur l'infrastructure}

Oui, surtout pour la modification dynamique. Si l'IA change le port de mon serveur de test d\`{e}s
qu'il est cibl\'{e} par un scan, cela me prot\`{e}ge sans interrompre la connexion avec mes coll\`{e}gues
de travaux pratiques.

\bigskip
\textbf{Q9. Limites des m\'{e}thodes actuelles}

Cela exige une surveillance constante du trafic. Il m'est impossible d'avancer sur mon code et
de surveiller le r\'{e}seau simultan\'{e}ment --- c'est une perte de temps significative.

\bigskip
\textbf{Q10. Caract\'{e}ristiques souhait\'{e}es d'une solution IA}

J'aimerais un tableau de bord web tr\`{e}s l\'{e}ger indiquant simplement quels ports sont
actuellement cibl\'{e}s par des requ\^{e}tes externes et le niveau de menace global.

\bigskip
\textbf{Q11. Craintes avant d'autoriser une IA \`{a} interagir avec le r\'{e}seau}

Ma principale crainte est que ce mod\`{e}le ralentisse les performances de ma propre machine
s'il consomme trop de ressources processeur pour analyser les flux en continu.

\bigskip
\textbf{Q12. Bilan et disposition \`{a} adopter la solution}

Si ce syst\`{e}me intelligent et autonome, garantissant un faible taux d'erreur, \'{e}tait d\'{e}ploy\'{e}, je
l'adopterais sans h\'{e}sitation. Cela s\'{e}curiserait mes projets locaux sans effort suppl\'{e}mentaire
de ma part.

\clearpage

% ============================================================
%  ENTRETIEN 5 — FAHD
% ============================================================
\section{Entretien 5 --- Fahd}

\subsection*{Informations sur l'interlocuteur}

\begin{tabular}{@{}B{4cm} L{10.5cm}@{}}
  \hline
  Nom et pr\'{e}nom      & Fahd \\
  \hline
  Poste / Fonction   & \'{E}tudiant en 1$^{re}$ ann\'{e}e du cycle ing\'{e}nieur en informatique \\
  \hline
  Entreprise         & --- \\
  \hline
  Secteur            & Cybers\'{e}curit\'{e} (contexte acad\'{e}mique) \\
  \hline
  Date de l'entretien & 28 avril 2026 \\
  \hline
  Lieu / Format      & Microsoft Teams \\
  \hline
\end{tabular}

\subsection*{Questions et r\'{e}ponses}

\textbf{Q1. Pr\'{e}sentation et architecture r\'{e}seau utilis\'{e}e}

Je m'appelle Fahd. Je suis \'{e}tudiant en premi\`{e}re ann\'{e}e du cycle ing\'{e}nieur en informatique. Je
me connecte au r\'{e}seau de l'\'{e}cole avec mon PC personnel pour suivre les cours et r\'{e}aliser des
travaux pratiques qui n\'{e}cessitent souvent des communications multi-ports avec des machines
virtuelles.

\bigskip
\textbf{Q2. Principaux d\'{e}fis li\'{e}s \`{a} l'infrastructure r\'{e}seau}

Les d\'{e}connexions intempestives constituent le principal d\'{e}fi. Parfois, le r\'{e}seau me
d\'{e}connecte simplement parce que je lance un outil de diagnostic n\'{e}cessaire \`{a} un projet. Le
filtrage est tr\`{e}s restrictif.

\bigskip
\textbf{Q3. Exp\'{e}rience r\'{e}cente d'anomalie r\'{e}seau}

R\'{e}cemment, je lan\c{c}ais un script de test r\'{e}seau totalement inoffensif vers une machine virtuelle
et, soudainement, il m'\'{e}tait impossible de charger la moindre page web. Mon adresse IP avait
\'{e}t\'{e} mise sur liste noire par la passerelle de l'\'{e}cole.

\bigskip
\textbf{Q4. Causes probables des perturbations furtives}

Je pense que des r\`{e}gles de s\'{e}curit\'{e} trop rigides s'activent pour un rien. D\`{e}s que le trafic sort
du cadre web standard, le syst\`{e}me de d\'{e}tection r\'{e}agit de mani\`{e}re disproportionn\'{e}e.

\bigskip
\textbf{Q5. Diff\'{e}rencier une panne d'une attaque}

C'est difficile \`{a} d\'{e}terminer. En g\'{e}n\'{e}ral, je fais un ping vers la passerelle : si elle r\'{e}pond mais
que tout le reste est bloqu\'{e}, je comprends que j'ai \'{e}t\'{e} filtr\'{e} arbitrairement par le syst\`{e}me de
d\'{e}tection d'intrusion.

\bigskip
\textbf{Q6. D\'{e}tection des tentatives de reconnaissance}

En tant qu'utilisateur, je ne les d\'{e}tecte pas du tout. Je subis uniquement les cons\'{e}quences
lorsque l'administrateur r\'{e}seau d\'{e}cide de bloquer des segments entiers par pr\'{e}caution apr\`{e}s
avoir d\'{e}tect\'{e} un scan.

\bigskip
\textbf{Q6b. Outils et proc\'{e}dures utilis\'{e}s}

Je n'ai pas les droits d'administration sur le mat\'{e}riel de l'\'{e}cole, donc je ne peux rien faire \`{a}
part red\'{e}marrer ma machine ou changer de connexion en esp\'{e}rant que le blocage soit lev\'{e}.

\bigskip
\textbf{Q7. Int\'{e}gration de l'IA pour d\'{e}tecter les attaques en temps r\'{e}el}

Si l'IA est correctement entra\^{i}n\'{e}e et qu'elle r\'{e}duit drastiquement les faux positifs en
distinguant un vrai scan d'un travail pratique \'{e}tudiant, je suis totalement pour.

\bigskip
\textbf{Q8. Pertinence de l'IA pour intervenir sur l'infrastructure}

Bloquer une IP automatiquement : oui, \`{a} condition qu'on puisse facilement contester le blocage
en cas d'erreur. L'id\'{e}e de modifier les acc\`{e}s dynamiquement me semble tr\`{e}s pertinente pour
esquiver un attaquant sans tout couper.

\bigskip
\textbf{Q9. Limites des m\'{e}thodes actuelles}

Les faux positifs, principalement. Les r\`{e}gles actuelles ne tiennent absolument pas compte du
contexte acad\'{e}mique de mes requ\^{e}tes et me bloquent sans raison valable, ce qui est
particuli\`{e}rement probl\'{e}matique lors des rendus de projets.

\bigskip
\textbf{Q10. Caract\'{e}ristiques souhait\'{e}es d'une solution IA}

Il faut avant tout de la transparence. Si mon flux est bloqu\'{e}, je veux une notification claire
expliquant exactement quelle caract\'{e}ristique de mon trafic a conduit l'IA \`{a} prendre cette
d\'{e}cision.

\bigskip
\textbf{Q11. Craintes avant d'autoriser une IA \`{a} interagir avec le r\'{e}seau}

Le taux de faux positifs doit \^{e}tre tr\`{e}s faible. Je ne veux plus \^{e}tre emp\^{e}ch\'{e} de rendre un
projet \`{a} cause d'une r\`{e}gle de s\'{e}curit\'{e} mal calibr\'{e}e.

\bigskip
\textbf{Q12. Bilan et disposition \`{a} adopter la solution}

Si cela met fin aux blocages injustifi\'{e}s que nous subissons actuellement, je serais le premier
\`{a} plaider pour son installation.

\clearpage

% ============================================================
%  ANNEXE — GUIDE D'ENTRETIEN
% ============================================================
\section*{Annexe --- Guide d'entretien}
\addcontentsline{toc}{section}{Annexe --- Guide d'entretien}

\begin{center}
  \textit{IA \& Automatisation de la d\'{e}tection d'intrusions --- Analyse des vuln\'{e}rabilit\'{e}s r\'{e}seau}
\end{center}

\begin{tabular}{@{}B{4cm} L{10.5cm}@{}}
  \hline
  Module         & Intelligence Artificielle \\
  \hline
  Statut         & Projet acad\'{e}mique \\
  \hline
  \'{E}tudiant 1 & El Kartouti Anas \\
  \hline
  \'{E}tudiant 2 & Lekhbioui Adil \\
  \hline
  \'{E}tudiant 3 & Moulay Anas \\
  \hline
  \'{E}tudiant 4 & Hammoubel El Yazid \\
  \hline
  \'{E}tudiant 5 & Nafia Khadija \\
  \hline
\end{tabular}

\bigskip
\subsection*{1. Phase d'introduction (mise en situation)}

\textbf{Q1.1 :} \og{} Pourriez-vous vous pr\'{e}senter bri\`{e}vement et d\'{e}crire votre r\^{o}le ainsi que le
type d'architecture r\'{e}seau (LAN, Cloud, r\'{e}seaux hybrides) que vous g\'{e}rez ou utilisez au
quotidien~? \fg{}

\bigskip

\textbf{Q1.2 :} \og{} De mani\`{e}re g\'{e}n\'{e}rale, quels sont les d\'{e}fis ou les incidents majeurs auxquels
vous \^{e}tes le plus fr\'{e}quemment confront\'{e} pour garantir la disponibilit\'{e} et la s\'{e}curit\'{e} de cette
infrastructure~? \fg{}

\subsection*{2. Phase de d\'{e}veloppement (exploration du probl\`{e}me cible)}

\textbf{Q2.1 :} \og{} Pourriez-vous me d\'{e}crire en d\'{e}tail une exp\'{e}rience r\'{e}cente o\`{u} vous avez
constat\'{e} un comportement suspect, une anomalie de trafic ou une baisse de performance
inexpliquée sur le r\'{e}seau~? \fg{}

\bigskip

\textbf{Q2.2 :} \og{} Selon votre exp\'{e}rience, lorsque le r\'{e}seau subit ce type de perturbations de
mani\`{e}re furtive, quelles en sont les causes les plus courantes~? \fg{}

\bigskip

\textbf{Q2.3 :} \og{} Face \`{a} ces anomalies, comment faites-vous la diff\'{e}rence entre une simple panne
de composant, une inondation de trafic malveillante (DDoS/surcharge) ou une infection par
un malware interne~? \fg{}

\bigskip

\textbf{Q2.4 :} \og{} Les attaquants proc\`{e}dent souvent par des phases de reconnaissance avant de
frapper. Comment percevez-vous ou d\'{e}tectez-vous ces tentatives de cartographie silencieuse
(comme l'exploration de vos services ou les scans de ports) sur votre p\'{e}rim\`{e}tre~? \fg{}

\subsection*{3. Phase d'approfondissement (limites actuelles et transition vers l'innovation)}

\textbf{Q3.1 :} \og{} Actuellement, quelles proc\'{e}dures, journaux ou outils mobilisez-vous pour
isoler la cause d'une intrusion ou d'un scan, et comment proc\'{e}dez-vous pour bloquer la
menace~? \fg{}

\bigskip

\textbf{Q3.2 :} \og{} Quelles sont les limites techniques ou humaines que vous rencontrez avec ces
m\'{e}thodes actuelles (temps de r\'{e}action, faux positifs, charge de travail)~? \fg{}

\bigskip

\textbf{Q3.3 :} \og{} Face \`{a} la complexit\'{e} des attaques, que penseriez-vous de l'int\'{e}gration d'un
mod\`{e}le d'intelligence artificielle capable d'analyser le trafic en temps r\'{e}el pour d\'{e}tecter de
mani\`{e}re autonome ces comportements suspects~? \fg{}

\bigskip

\textbf{Q3.4 :} \og{} Pensez-vous qu'un mod\`{e}le IA serait pertinent non seulement pour d\'{e}tecter,
mais aussi pour intervenir directement sur l'infrastructure (ex.~: bloquer une IP, modifier
dynamiquement l'acc\`{e}s) afin de vous faire gagner du temps~? \fg{}

\subsection*{4. Phase de conclusion (crit\`{e}res d'acceptabilit\'{e} de la solution)}

\textbf{Q4.1 :} \og{} Pour qu'une solution bas\'{e}e sur l'IA vous soit r\'{e}ellement utile, quelles
devraient \^{e}tre ses caract\'{e}ristiques principales (ex.~: type de tableau de bord, alertes en
temps r\'{e}el, visibilit\'{e} du r\'{e}seau)~? \fg{}

\bigskip

\textbf{Q4.2 :} \og{} Quelles seraient vos exigences strictes ou vos craintes (ex.~: confidentialit\'{e}
des flux, blocage accidentel d'un trafic l\'{e}gitime) avant d'autoriser une IA \`{a} interagir avec vos
routeurs/pare-feux~? \fg{}

\bigskip

\textbf{Q4.3 :} \og{} En r\'{e}sum\'{e}, si ce syst\`{e}me intelligent et autonome, garantissant un faible taux
d'erreur, \'{e}tait d\'{e}ploy\'{e} sur votre infrastructure, seriez-vous pr\^{e}t \`{a} l'adopter dans vos
op\'{e}rations quotidiennes~? \fg{}

\clearpage

% ============================================================
%  CONCLUSION
% ============================================================
\section*{Conclusion}
\addcontentsline{toc}{section}{Conclusion}

Ces entretiens nous ont permis de d\'{e}passer le cadre purement acad\'{e}mique et d'appr\'{e}hender
la d\'{e}tection des attaques r\'{e}seau sous un angle professionnel concret. Les \'{e}changes avec les
experts ont confirm\'{e} la pertinence de notre approche bas\'{e}e sur le machine learning, tout en
nous sensibilisant aux d\'{e}fis r\'{e}els li\'{e}s \`{a} la qualit\'{e} des donn\'{e}es, aux faux positifs et \`{a}
l'\'{e}volution constante des techniques d'attaque.

En tant qu'\'{e}tudiants en troisi\`{e}me ann\'{e}e du cycle ing\'{e}nieur, sp\'{e}cialit\'{e} R\'{e}seaux et Syst\`{e}mes
de T\'{e}l\'{e}communications \`{a} l'ENSA K\'{e}nitra, ces rencontres ont \'{e}t\'{e} particuli\`{e}rement
enrichissantes. Elles nous ont donn\'{e} une vision plus claire des comp\'{e}tences attendues dans
le domaine et nous ont motiv\'{e}s \`{a} approfondir notre ma\^{i}trise des outils de cybers\'{e}curit\'{e} et
d'analyse de trafic r\'{e}seau.

Nous tenons \`{a} remercier chaleureusement tous les professionnels qui ont accept\'{e} de nous
consacrer de leur temps et de partager leur expertise avec nous.

\vspace{2em}
\noindent\rule{\linewidth}{0.4pt}
\begin{center}
  \small Compte rendu des entretiens ---
  Module Intelligence Artificielle ---
  ENSA K\'{e}nitra --- 2025 / 2026
\end{center}

\end{document}
